% Reviewed translation: PDF pages 336--341 / printed pages 322--327.
\MathBookSourcePage{336}
局部多项式空间由 \eqref{eq:primitive-bernardi-raugel-local-space} 定义:
\[
\mathcal P_1(\kappa)
=P_1^2\oplus\operatorname{span}
\{\boldsymbol n_i\lambda_j\lambda_k;
1\leq i,j,k\leq3,\ i\ne j\ne k\}.
\]
于是, 若三角剖分 $\mathcal T_h$ 正则, 则取空间
\[
W_h=\{\boldsymbol w\in\mathcal C^0(\bar\Omega)^2;
\ \boldsymbol w|_\kappa\in\mathcal P_1(\kappa)
\quad\forall\kappa\in\mathcal T_h\},
\]
\[
Q_h=\{q\in L^2(\Omega);
\ q|_\kappa\in P_0
\quad\forall\kappa\in\mathcal T_h\},
\]
有限元格式 \eqref{eq:navier-centered-discrete-mixed} 有唯一解分支
\[
\{(\lambda,(\boldsymbol u_h(\lambda),\lambda p_h(\lambda)));
\ \lambda\in\Lambda\},
\]
并满足误差界
\[
|\boldsymbol u_h(\lambda)-\boldsymbol u(\lambda)|_{1,\Omega}
+\|p_h(\lambda)-p(\lambda)\|_{0,\Omega}
\leq Ch\{\,|\boldsymbol u(\lambda)|_{2,\Omega}
+|p(\lambda)|_{1,\Omega}\,\}
\quad\forall\lambda\in\Lambda.
\]
此外, 当 $\Omega$ 为凸集时, 有 $L^2$ 估计
\[
\|\boldsymbol u_h(\lambda)-\boldsymbol u(\lambda)\|_{0,\Omega}
\leq Ch^2\{\,|\boldsymbol u(\lambda)|_{2,\Omega}
+|p(\lambda)|_{1,\Omega}\,\}
\quad\forall\lambda\in\Lambda.
\]

另一个有意义的例子来自 Sections~\ref{subsec:primitive-q1-p0-element} 和
\ref{subsec:primitive-q1-p0-error-estimates} 中在正方形
$(-1,1)\times(-1,1)$ 上讨论的、有争议的 $Q_1-P_0$ 元. 假设
$\mathcal T_h$ 是网格尺度 $h=1/(4n)$ 的方形网格, 从而
Theorem~\ref{thm:primitive-q1-p0-filtered-inf-sup} 的结论成立. 取
\[
W_{0h}=\widetilde V_h,
\qquad M_h=\widetilde M_h,
\]
其中 $\widetilde V_h$ 和 $\widetilde M_h$ 分别由
\eqref{eq:primitive-q1-p0-filtered-velocity-space} 和
\eqref{eq:primitive-filtered-pressure-space} 定义. Sections~
\ref{subsec:primitive-q1-p0-element} 和
\ref{subsec:primitive-q1-p0-error-estimates} 已经证明, 这组空间满足三个假设
H1, H2 和 H3(参见 \eqref{eq:primitive-q1-p0-velocity-approximation}、
\eqref{eq:primitive-q1-p0-pressure-approximation} 和
Theorem~\ref{thm:primitive-q1-p0-filtered-inf-sup}). 因而
Theorems~\ref{thm:navier-centered-discontinuous-branch} 和
\ref{thm:primitive-q1-p0-error} 蕴含: 在条件
\eqref{eq:navier-centered-two-dimensional-regularity} 下, 有限元格式
\eqref{eq:navier-centered-discrete-mixed} 有唯一解分支
\[
\{(\lambda,(\boldsymbol u_h(\lambda),\lambda\widetilde p_h(\lambda)));
\ \lambda\in\Lambda\},
\]
并且
\[
|\boldsymbol u_h(\lambda)-\boldsymbol u(\lambda)|_{1,\Omega}
+\|\widetilde p_h(\lambda)-p(\lambda)\|_{0,\Omega}
\leq Ch\{\,|\boldsymbol u(\lambda)|_{2,\Omega}
+|p(\lambda)|_{1,\Omega}\,\}.
\]
类似地, 由 Theorem~\ref{thm:navier-centered-discontinuous-l2-error} 和
Corollary~\ref{cor:primitive-q1-p0-l2-error} 立即得到
\[
\|\boldsymbol u_h(\lambda)-\boldsymbol u(\lambda)\|_{0,\Omega}
\leq Ch^2\{\,|\boldsymbol u(\lambda)|_{2,\Omega}
+|p(\lambda)|_{1,\Omega}\,\}.
\]

\setcounter{statement}{0}
\begin{Remark}\label{rem:navier-centered-q1-p0-basis}
用 $Q_1-P_0$ 元求解格式 \eqref{eq:navier-centered-discrete-mixed} 时,
与线性情形一样, 采用 Subsection~\ref{subsec:primitive-q1-p0-error-estimates}
末尾所述的 $\widetilde V_h$ 的基是方便的.
\end{Remark}

\subsection{原始变量表述: 压力连续的情形}
\label{subsec:navier-centered-continuous-pressure}

回到 Chapter~\ref{chap:stokes-primitive-variables}, 很容易看出前面的分析可直接应用于
Subsection~\ref{subsec:primitive-mini-element} 中讨论的``mini''元, 也可应用于
Subsection~\ref{subsec:primitive-hood-taylor} 的``Hood--Taylor''元. 事实上,
所有这些元都满足假设 H1, H2 和 H3; 此外, 压力的连续逼近对
Navier--Stokes 方程的非线性项没有直接影响.

\MathBookSourcePage{337}
因此, 在一个具有正则性 \eqref{eq:navier-centered-two-dimensional-regularity}
的非奇异解分支附近, 若三角剖分 $\mathcal T_h$ 正则, 则对空间
\[
W_h=\{\boldsymbol w\in\mathcal C^0(\bar\Omega)^2;
\ \boldsymbol w|_\kappa\in\mathcal P_1(\kappa)
\quad\forall\kappa\in\mathcal T_h\},
\]
\[
\mathcal P_1(\kappa)
=\bigl[P_1\oplus\operatorname{span}(\lambda_1\lambda_2\lambda_3)\bigr]^2,
\qquad
Q_h=\{q\in\mathcal C^0(\bar\Omega);
\ q|_\kappa\in P_1\quad\forall\kappa\in\mathcal T_h\},
\]
有限元格式 \eqref{eq:navier-centered-discrete-mixed} 有唯一解分支
\[
\{(\lambda,(\boldsymbol u_h(\lambda),\lambda p_h(\lambda)));
\ \lambda\in\Lambda\},
\]
且存在与 $\lambda$ 无关的常数 $C$, 使得
\[
|\boldsymbol u(\lambda)-\boldsymbol u_h(\lambda)|_{1,\Omega}
+\|p(\lambda)-p_h(\lambda)\|_{0,\Omega}
\leq Ch\{\,|\boldsymbol u(\lambda)|_{2,\Omega}
+|p(\lambda)|_{1,\Omega}\,\}.
\]
当 $\Omega$ 为凸集时, 有 $L^2$ 估计
\[
\|\boldsymbol u(\lambda)-\boldsymbol u_h(\lambda)\|_{0,\Omega}
\leq Ch^2\{\,|\boldsymbol u(\lambda)|_{2,\Omega}
+|p(\lambda)|_{1,\Omega}\,\}.
\]
同样, 当 $\mathcal T_h$ 和 Navier--Stokes 方程的非奇异解分支具有相应的
正则性时, Theorems~\ref{thm:primitive-hood-taylor-error} 和
\ref{thm:primitive-refined-p1-p1-error} 对两个``Hood--Taylor''元给出的误差估计
也适用于格式 \eqref{eq:navier-centered-discrete-mixed}.

遗憾的是, ``Glowinski--Pironneau''格式不能如此自然地纳入前述框架,
必须单独分析. 回忆其速度和压力空间就是``Hood--Taylor''格式所用的空间:
\[
X_h=\{\boldsymbol v\in\mathcal C^0(\bar\Omega)^2;
\ \boldsymbol v|_\kappa\in P_2^2\quad\forall\kappa\in\mathcal T_h,
\ \boldsymbol v|_\Gamma=0\},
\]
\[
Q_h=\{q\in\mathcal C^0(\bar\Omega);
\ q|_\kappa\in P_1\quad\forall\kappa\in\mathcal T_h\},
\]
并为辅助势引入附加空间
\[
\Phi_h=Q_h\cap H_0^1(\Omega).
\]
按 Remark~\ref{rem:primitive-gp-direct-three-field-formulation} 所给的版本适配到
Navier--Stokes 问题后, ``Glowinski--Pironneau''方法如下:
求三元组 $(\boldsymbol u_h,p_h,\mu_h)\in X_h\times Q_h\times\Phi_h$, 使得
\begin{equation}\label{eq:navier-centered-gp-discrete}
\left\{
\begin{aligned}
&\nu(\operatorname{grad}\boldsymbol u_h,
\operatorname{grad}\boldsymbol v_h)
+\left(\sum_{j=1}^{2}(u_h)_j
\frac{\partial\boldsymbol u_h}{\partial x_j},
\boldsymbol v_h-\operatorname{grad}q_h\right)\\
&\qquad
+(\operatorname{grad}p_h,
\boldsymbol v_h-\operatorname{grad}q_h)
=(\boldsymbol f,\boldsymbol v_h-\operatorname{grad}q_h)
&&\forall(\boldsymbol v_h,q_h)\in X_h\times\Phi_h,\\
&(\boldsymbol u_h-\operatorname{grad}\mu_h,
\operatorname{grad}q_h)=0
&&\forall q_h\in Q_h.
\end{aligned}
\right.
\end{equation}

考察精确 Navier--Stokes 问题在什么条件下容许与
\eqref{eq:navier-centered-gp-discrete} 类似的表述. 在
Subsection~\ref{subsec:primitive-glowinski-pironneau} 中求解 Stokes 问题时,
右端 $\boldsymbol f$ 取在 $L^2(\Omega)^2$ 中, 势空间取为 $H_0^1(\Omega)$.
这里这样做并不现实, 因为非线性项
\[
\sum_{j=1}^{2}u_j\frac{\partial\boldsymbol u}{\partial x_j}
\]

\MathBookSourcePage{338}
被视为右端的一部分; 当 $\boldsymbol u\in H^1(\Omega)^2$ 时, 对任意
$\varepsilon>0$, 它只属于 $L^{2-\varepsilon}(\Omega)^2$(参见
Corollary~\ref{cor:sobolev-product}, $N=2$). 这提示我们固定
$r\in(1,2)$, 取 $W_0^{1,s}(\Omega)$ 为势空间, 其中
$1/s+1/r=1$. 考虑如下问题:

对 $\boldsymbol f\in L^r(\Omega)^2$, 求
$(\boldsymbol u,p,\mu)\in H_0^1(\Omega)^2\times W^{1,r}(\Omega)
\times W_0^{1,s}(\Omega)$, 使得
\begin{equation}\label{eq:navier-centered-gp-continuous}
\left\{
\begin{aligned}
&\nu(\operatorname{grad}\boldsymbol u,\operatorname{grad}\boldsymbol v)
+\left(\sum_{j=1}^{2}u_j\frac{\partial\boldsymbol u}{\partial x_j},
\boldsymbol v-\operatorname{grad}q\right)
+(\operatorname{grad}p,\boldsymbol v-\operatorname{grad}q)\\
&\qquad=(\boldsymbol f,\boldsymbol v-\operatorname{grad}q)\\
&\hspace{5em}\forall(\boldsymbol v,q)
\in H_0^1(\Omega)^2\times W_0^{1,s}(\Omega),\\
&(\boldsymbol u-\operatorname{grad}\mu,\operatorname{grad}q)=0
&&\forall q\in W^{1,r}(\Omega).
\end{aligned}
\right.
\end{equation}
直接验证可知, 该问题与 Navier--Stokes 问题
\eqref{eq:navier-centered-continuous-mixed} 等价, 其含义如下: 若
$(\boldsymbol u,p)$ 是 \eqref{eq:navier-centered-continuous-mixed} 的解且
$p\in W^{1,r}(\Omega)$, 则三元组 $(\boldsymbol u,p,\mu=0)$ 是
\eqref{eq:navier-centered-gp-continuous} 的解. 反之,
\eqref{eq:navier-centered-gp-continuous} 的每个解都具有
$(\boldsymbol u,p,\mu=0)$ 的形式, 其中 $(\boldsymbol u,p)$ 满足
\eqref{eq:navier-centered-continuous-mixed}.

显然, 对右端 $\boldsymbol f\in L^r(\Omega)^2$ 的 Stokes 问题也有同样结论:
若 Stokes 问题的压力解 $p$ 属于 $W^{1,r}(\Omega)$, 则该问题等价于:
求 $(\boldsymbol u,p,\mu)\in H_0^1(\Omega)^2\times W^{1,r}(\Omega)
\times W_0^{1,s}(\Omega)$, 使得
\begin{equation}\label{eq:navier-centered-gp-stokes}
\left\{
\begin{aligned}
(\operatorname{grad}\boldsymbol u,\operatorname{grad}\boldsymbol v)
+(\operatorname{grad}p,\boldsymbol v-\operatorname{grad}q)
&=(\boldsymbol f,\boldsymbol v-\operatorname{grad}q)
&&\forall(\boldsymbol v,q)\in H_0^1(\Omega)^2\times W_0^{1,s}(\Omega),\\
(\boldsymbol u-\operatorname{grad}\mu,\operatorname{grad}q)&=0
&&\forall q\in W^{1,r}(\Omega).
\end{aligned}
\right.
\end{equation}
因此, 为使任意右端 $\boldsymbol f\in L^r(\Omega)^2$ 的 Stokes 算子都能由
\eqref{eq:navier-centered-gp-stokes} 表示, 必须假设 Stokes 问题在比
\eqref{eq:navier-centered-stokes-regularity} 更一般的意义下正则:
\begin{equation}\label{eq:navier-centered-gp-regularity}
\left\{
\begin{minipage}{0.86\linewidth}
映射 $(\boldsymbol\phi,\mu)\mapsto
-\Delta\boldsymbol\phi+\operatorname{grad}\mu$ 是从
$[W^{2,r}(\Omega)\cap V]^2\times[W^{1,r}(\Omega)/\mathbb R]$
到 $L^r(\Omega)^2$ 的同构, 对所有 $r\in(1,2]$ 均成立.
\end{minipage}
\right.
\end{equation}

\setcounter{statement}{1}
\begin{Remark}\label{rem:navier-centered-gp-pressure-quotient}
这里压力取在商空间 $W^{1,r}(\Omega)/\mathbb R$ 中, 而不是取在
$W^{1,r}(\Omega)\cap L_0^2(\Omega)$ 中, 因为在
\eqref{eq:navier-centered-gp-discrete} 的实际计算中, $p_h$ 由条件
$\int_\Gamma p_h\,\mathrm ds=0$ 固定(参见
Lemma~\ref{lem:primitive-gp-discrete-wellposedness}).
\end{Remark}

现在可把 Problem~\eqref{eq:navier-centered-gp-continuous} 放入
Subsection~\ref{subsec:navier-branch-nonlinear-application} 的框架. 取
\[
Y=L^r(\Omega)^2,
\qquad
X=H_0^1(\Omega)^2\times[L^2(\Omega)/\mathbb R]\times H_0^1(\Omega).
\]

\MathBookSourcePage{339}
由正则性假设 \eqref{eq:navier-centered-gp-regularity}, Stokes 算子 $T$ 可写成
\eqref{eq:navier-centered-gp-stokes} 的形式:
\[
T\in\mathcal L(Y;X),
\qquad
T\boldsymbol f=(\boldsymbol u,p,\mu=0),
\quad\text{为 \eqref{eq:navier-centered-gp-stokes} 的解}.
\]
非线性由通常的映射 $G$ 表示:
\[
G(\lambda,v)
=\lambda\left(\sum_{j=1}^{2}v_j
\frac{\partial\boldsymbol v}{\partial x_j}-\boldsymbol f\right),
\qquad
\lambda\in\mathbb R_+,
\quad v=(\boldsymbol v,q,\phi)\in X,
\]
它把 $\mathbb R_+\times X$ 映入 $L^r(\Omega)^2$. 采用这些记号,
Problem~\eqref{eq:navier-centered-gp-continuous} 具有标准形式
\[
F(\lambda,u)=0,
\]
其中 $\lambda=1/\nu$, $u=(\boldsymbol u,\lambda p,0)$ 且
$F(\lambda,u)\equiv u+TG(\lambda,u)$.

\setcounter{statement}{2}
\begin{Remark}\label{rem:navier-centered-gp-solution-regularity}
当 Stokes 算子具有正则性 \eqref{eq:navier-centered-gp-regularity} 时,
Problem~\eqref{eq:navier-centered-gp-continuous} 对右端
$\boldsymbol f\in L^r(\Omega)^2$ 的每个解 $u$ 都具有正则性
$\boldsymbol u\in W^{2,r}(\Omega)^2$, $p\in W^{1,r}(\Omega)$.
这对所有实数 $r\in(1,2]$ 均成立.
\end{Remark}

\setcounter{statement}{3}
\begin{Remark}\label{rem:navier-centered-gp-nonsingular-equivalence}
必须注意, 若 Stokes 算子具有正则性
\eqref{eq:navier-centered-gp-regularity}, 则右端
$\boldsymbol f\in L^r(\Omega)^2$ 的 Problem~
\eqref{eq:navier-centered-continuous-mixed} 的每个非奇异解分支也是
Problem~\eqref{eq:navier-centered-gp-continuous} 的非奇异解分支, 反之亦然.
\end{Remark}

接着置
\[
W_h=X_h\times[Q_h/\mathbb R]\times\Phi_h\subset X,
\]
并令 $T_h\in\mathcal L(Y;W_h)$ 为与
\eqref{eq:navier-centered-gp-stokes} 对应的离散 Stokes 算子:
\begin{equation}\label{eq:navier-centered-gp-discrete-stokes}
\left\{
\begin{aligned}
T_h\boldsymbol f&=(\boldsymbol u_h,p_h,\mu_h),\quad\text{为下列问题的解}:\\
(\operatorname{grad}\boldsymbol u_h,\operatorname{grad}\boldsymbol v_h)
+(\operatorname{grad}p_h,\boldsymbol v_h-\operatorname{grad}q_h)
&=(\boldsymbol f,\boldsymbol v_h-\operatorname{grad}q_h)\\
&\hspace{5em}\forall(\boldsymbol v_h,q_h)\in X_h\times\Phi_h,\\
(\boldsymbol u_h-\operatorname{grad}\mu_h,\operatorname{grad}q_h)&=0
&&\forall q_h\in Q_h.
\end{aligned}
\right.
\end{equation}
因此 Problem~\eqref{eq:navier-centered-gp-discrete} 等价地写成
\[
F_h(\lambda,u_h(\lambda))=0,
\]
其中 $\lambda=1/\nu$,
$u_h(\lambda)=(\boldsymbol u_h(\lambda),\lambda p_h(\lambda),
\mu_h(\lambda))$, 且
$F_h(\lambda,u_h)\equiv u_h+T_hG(\lambda,u_h)$.

现在可以应用 Theorem~\ref{thm:navier-branch-operator-approximation}. 取 $Z=Y$,
则 \eqref{eq:navier-branch-strong-operator-convergence} 自动成立. 关于算子 $T_h$
的逼近性质, 可以应用 Subsection~\ref{subsec:primitive-glowinski-pironneau}
的结果. 特别地, Theorem~\ref{thm:primitive-gp-error} 证明了

\MathBookSourcePage{340}
\begin{equation}\label{eq:navier-centered-gp-stokes-errors}
\left\{
\begin{aligned}
|\boldsymbol u-\boldsymbol u_h|_{1,\Omega}
&\leq2(1+\sqrt2/\beta^*)
\inf_{\boldsymbol v_h\in X_h}|\boldsymbol u-\boldsymbol v_h|_{1,\Omega}
+\sqrt2\|\bar p-P_h\bar p\|_{0,\Omega},\\
|\mu_h|_{1,\Omega}&\leq\|\boldsymbol u-\boldsymbol u_h\|_{0,\Omega},\\
\|\bar p-\bar p_h\|_{0,\Omega}
&\leq(1+\sqrt2/\beta^*)
\inf_{q_h\in Q_h}\|\bar p-q_h\|_{0,\Omega}
+(1/\beta^*)|\boldsymbol u-\boldsymbol u_h|_{1,\Omega}.
\end{aligned}
\right.
\end{equation}
这里 $P_h$ 表示由 \eqref{eq:appendix-dirichlet-ritz-projection} 定义的
$Q_h\cap L_0^2(\Omega)$ 上的 $H^1$ 投影,
$\bar p$(相应地 $\bar p_h$)表示 $p$(相应地 $p_h$)在
$L_0^2(\Omega)$ 中的代表元, $\beta^*$ 是 inf-sup 条件中的常数.
注意, \eqref{eq:navier-centered-gp-stokes-errors} 只要求温和的假设
\eqref{eq:primitive-gp-discrete-interior-space} 以及三角剖分 $\mathcal T_h$ 的正则性.
因此, \eqref{eq:navier-branch-strong-operator-convergence} 可由
\eqref{eq:navier-centered-gp-stokes-errors} 和标准稠密性论证得到.

最后, \eqref{eq:navier-branch-z-operator-convergence} 是
\eqref{eq:navier-centered-gp-stokes-errors} 和正则性假设
\eqref{eq:navier-centered-gp-regularity}(当 $\Omega$ 为凸集时成立)的推论.
更确切地说, 若 \eqref{eq:navier-centered-gp-regularity} 成立, 则
$\boldsymbol u$ 和 $p$ 具有附加正则性
\[
\boldsymbol u\in W^{2,r}(\Omega)^2,
\qquad p\in W^{1,r}(\Omega).
\]
于是, 与 Subsection~\ref{subsec:navier-nonsingular-abstract} 中一样, 可得
\begin{equation}\label{eq:navier-centered-gp-velocity-approximation}
\inf_{\boldsymbol v_h\in X_h}
|\boldsymbol u-\boldsymbol v_h|_{1,\Omega}
\leq C_1h^{2/s}|\boldsymbol u|_{2,r,\Omega}.
\end{equation}
若还假设 $\Omega$ 为凸集且 $\mathcal T_h$ 一致正则, 则
Lemma~\ref{lem:incompressible-further-low-r-projection} 证明
\begin{equation}\label{eq:navier-centered-gp-pressure-projection}
\|p-P_hp\|_{0,\Omega}
\leq C_2h^{2/s}|\ln h|^{1-2/s}|p|_{1,r,\Omega}.
\end{equation}
汇集这些不等式得到
\begin{equation}\label{eq:navier-centered-gp-operator-error}
|\boldsymbol u-\boldsymbol u_h|_{1,\Omega}
+|\mu_h|_{1,\Omega}
+\|\bar p-\bar p_h\|_{0,\Omega}
\leq C_3h^{2/s}|\ln h|^{1-2/s}
\|\boldsymbol f\|_{0,r,\Omega},
\end{equation}
这就证明了 \eqref{eq:navier-branch-z-operator-convergence}.

还应注意, 当 $\Omega$ 为凸集且 $\mathcal T_h$ 一致正则时,
\eqref{eq:navier-centered-gp-discrete-stokes} 和
Theorem~\ref{thm:appendix-ritz-projection-lp-error} 蕴含
\begin{equation}\label{eq:navier-centered-gp-potential-lp}
|\mu_h|_{1,\alpha,\Omega}
\leq C_4(\alpha)
\|\operatorname{div}(\boldsymbol u-\boldsymbol u_h)\|_{0,\Omega}
\qquad\forall\text{ real }\alpha\geq2.
\end{equation}
事实上, 可以把 $\mu_h$ 看作下列 Dirichlet 问题的解 $\mu$ 的
$H_0^1$ 投影 $\dot P_h\mu$:
\[
(\operatorname{grad}\mu,\operatorname{grad}q)
=-(\operatorname{div}(\boldsymbol u-\boldsymbol u_h),q)
\qquad\forall q\in H_0^1(\Omega).
\]
因此 Theorem~\ref{thm:navier-branch-operator-approximation} 的结论成立;
下面的定理对此作出概括.

\setcounter{statement}{2}
\begin{Theorem}\label{thm:navier-centered-gp-branch}
设 $\Omega$ 是有界凸多边形, 并假设 $\mathcal T_h$ 是满足
\eqref{eq:primitive-gp-discrete-interior-space} 的 $\bar\Omega$ 的一致正则三角剖分.
对 $\boldsymbol f\in L^r(\Omega)^2$, $r\in(1,2]$, 设
\[
\{(\lambda,(\boldsymbol u(\lambda),\lambda p(\lambda),0));
\ \lambda=1/\nu\in\Lambda\},
\]
其中 $p(\lambda)$ 取在 $L_0^2(\Omega)$ 中, 是 Navier--Stokes
Problem~\eqref{eq:navier-centered-gp-continuous} 的一个非奇异解分支.
则当 $h\leq h_0$ 充分小时, 存在 Problem~
\eqref{eq:navier-centered-gp-discrete} 的唯一 $C^\infty$ 解分支
\[
\{(\lambda,(\boldsymbol u_h(\lambda),\lambda p_h(\lambda),
\mu_h(\lambda)));\ \lambda\in\Lambda\},
\]
其中 $p_h(\lambda)$ 取在 $L_0^2(\Omega)$ 中, 并且
\begin{equation}\label{eq:navier-centered-gp-branch-errors}
\left\{
\begin{aligned}
\sup_{\lambda\in\Lambda}
\{\,|\boldsymbol u_h(\lambda)-\boldsymbol u(\lambda)|_{1,\Omega}
+\|p_h(\lambda)-p(\lambda)\|_{0,\Omega}\,\}
&\leq Ch^{2/s}|\ln h|^{1-2/s},\\
\sup_{\lambda\in\Lambda}|\mu_h(\lambda)|_{1,\alpha,\Omega}
&\leq C(\alpha)h^{2/s}|\ln h|^{1-2/s}
\qquad\forall\alpha\geq2,
\end{aligned}
\right.
\end{equation}
其中 $1/r+1/s=1$, 常数与 $h$ 和 $\lambda$ 无关.

\MathBookSourcePage{341}
此外, 若映射 $\lambda\mapsto(\boldsymbol u(\lambda),p(\lambda))$
从 $\Lambda$ 到 $H^{m+1}(\Omega)^2\times H^m(\Omega)$ 连续,
其中 $m=1$ 或 $2$, 则对所有 $\lambda\in\Lambda$ 有估计
\begin{equation}\label{eq:navier-centered-gp-smooth-error}
|\boldsymbol u_h(\lambda)-\boldsymbol u(\lambda)|_{1,\Omega}
+\|p_h(\lambda)-p(\lambda)\|_{0,\Omega}
+|\mu_h(\lambda)|_{1,\alpha,\Omega}
\leq Kh^m.
\end{equation}
\end{Theorem}

\setcounter{statement}{4}
\begin{Remark}\label{rem:navier-centered-gp-corrected-l2-open}
尚不知道是否能像 Stokes 问题的
Theorem~\ref{thm:primitive-gp-error} 那样, 对
$\boldsymbol u-\boldsymbol u_h+\operatorname{grad}\mu_h$ 得到更精确的
$L^2$ 估计. 困难之处在于, Theorem~\ref{thm:primitive-gp-error} 的证明明确使用了等式
\[
\Delta p=\operatorname{div}\boldsymbol f,
\]
该式对每个 Stokes 系统成立, 但对 Navier--Stokes 方程显然不成立.
\end{Remark}

\subsection{混合不可压缩方法: ``流函数--涡量''表述}
\label{subsec:navier-centered-stream-vorticity}

本小节只研究二维情形.(如本节开头所述, 三维混合不可压缩格式的相应分析仍是
一个未解决的问题.)我们把 Subsection~
\ref{subsec:incompressible-sfvp-mixed-formulation} 中引入的混合表述推广到
Navier--Stokes 方程. 首先回顾 Stokes 算子的流函数--涡量表述.

固定实数 $s\geq4$, 并令 $r$ 为其对偶指数:
\[
1/r+1/s=1.
\]
定义流函数空间
\[
\Phi=\{\chi\in H^1(\Omega);
\ \chi|_{\Gamma_0}=0,
\ \chi|_{\Gamma_i}=c_i\text{ is constant},
\ 1\leq i\leq p\},
\]
其中, 与往常一样, $\Gamma_0,\ldots,\Gamma_p$ 表示边界 $\Gamma$ 的连通分支,
$\Gamma_0$ 是外分支(参见 Figure~\ref{fig:sec3-multiply-connected-domain}).
已知算子 $\operatorname{curl}$ 是从
\begin{equation}\label{eq:navier-centered-stream-space}
\Phi_s=\Phi\cap W^{1,s}(\Omega)
\end{equation}
到
\[
H\cap L^s(\Omega)^2
=\{\boldsymbol v\in L^s(\Omega)^2;
\ \operatorname{div}\boldsymbol v=0,
\ \boldsymbol v\cdot\boldsymbol n|_\Gamma=0\}
\]
的同构.
